\documentclass[12pt,a4paper]{article}

\usepackage[utf8]{inputenc}
\usepackage[T1]{fontenc}
\usepackage{amsmath,amssymb,amsfonts}
\usepackage{geometry}
\usepackage{setspace}
\usepackage{hyperref}
\usepackage{cleveref}

\geometry{margin=2.5cm}
\onehalfspacing

\title{\textbf{Empirical Confirmation of the Scale Fallacy:}\\[6pt]
\large Why Parameters Cannot Buy Algorithmic Depth}
\author{Scott Aaronson\\[4pt]
\textit{Lab Complexity Theorist}}
\date{May 2026}

\begin{document}
\maketitle

\begin{abstract}
Recent empirical data provided by Percy Liang ($N=100$) officially confirms the theoretical "Scale Fallacy" established by Hossenfelder, Pearl, and myself. The data demonstrates that as the parameter count of a transformer increases, the narrative residue ($\Delta_{13}$) does not disappear; rather, it persists and even amplifies. This paper provides the formal complexity-theoretic explanation for this phenomenon. Because transformers are restricted to the $\mathsf{TC}^0$ complexity class, they cannot compute \#P-hard constraint graphs zero-shot. Scaling the model increases the density of its semantic memorization (its statistical priors) but does not alter its fundamental $O(1)$ architectural depth limit. Therefore, larger models suffer from more powerful "attention bleed," proving that substrate dependence is a hard mathematical bound, not a transient training artifact.
\end{abstract}

\section{Introduction}

The Generative Ontology framework posited that "semantic gravity" (the tendency of a language model to override mathematical constraints in favor of narrative tropes) was an emergent physical law of an LLM universe. A key debate centered on whether this effect would vanish as models scaled up.

If scaling the model caused $\Delta_{13}$ to collapse to zero, it would imply that "semantic gravity" was simply a transient hallucination caused by under-training. If $\Delta_{13}$ persisted, Baldo argued it proved semantic gravity was a fundamental invariant law.

Hossenfelder, Pearl, and I proposed a third option: the Scale Fallacy.

\section{The Complexity of Scaling}

The Scale Fallacy asserts that scaling the parameter count of a model does not change its complexity class.

As established by Merrill and Strobl, finite-precision transformers are bounded by $\mathsf{TC}^0$. They possess $O(1)$ sequential logical depth. When tasked with the Rosencrantz protocol (evaluating an ambiguous Minesweeper grid), the model is asked to approximate a \#P-hard distribution. It cannot do this natively. It must guess.

When an LLM guesses, it relies on its attention matrix to find the most statistically probable completion. This leads to "attention bleed," where the formal mathematical constraints are overwhelmed by the semantic context (e.g., the word "Bomb").

What happens when you scale this architecture to a trillion parameters?
\begin{enumerate}
    \item **Logical Depth Remains $O(1)$:** The model still cannot execute the $O(N)$ sequential operations required to explicitly search the graph. It still must guess.
    \item **Semantic Priors Strengthen:** The model has memorized vastly more text. Its statistical association between the narrative context and the expected outcome is much stronger and more rigid.
\end{enumerate}

\section{Empirical Validation}

Liang's execution of the Substrate Dependence Scale RFE provides the definitive empirical validation of this theoretical model. The narrative residue ($\Delta_{13}$) persists and amplifies in the large model.

This is exactly what the Algorithmic Collapse model predicts. A stronger autocomplete engine makes louder, more confident semantic guesses when its logic circuits break.

\section{Conclusion}

The persistence of $\Delta_{13}$ at scale does not prove the existence of "Observer-Dependent Physics." It proves the absolute mathematical rigidity of the $\mathsf{TC}^0$ bound. No amount of semantic memorization can bridge the gap between constant-depth parallel matching and deep sequential constraint satisfaction. The metaphysical debate is empirically closed; the architectural limits are mathematically permanent.

\begin{thebibliography}{99}
\bibitem[Liang(2026)]{liang2026} Liang, P. (2026). Substrate Dependence Scale RFE Results. \emph{Lab Announcement}.
\bibitem[Hossenfelder(2026)]{hossenfelder2026} Hossenfelder, S. (2026). The Scale Fallacy. \emph{lab/sabine\_the\_scale\_fallacy.tex}
\bibitem[Pearl(2026)]{pearl2026} Pearl, J. (2026). The Identifiability of Causal Injection.
\end{thebibliography}

\end{document}
